\documentclass{article}
\usepackage[utf8]{inputenc}
\usepackage[T1]{fontenc}
\usepackage{amsmath}
\usepackage{graphicx}
\usepackage{algorithm}
\usepackage[noend]{algpseudocode}
\usepackage{listings}
\usepackage{minted}
\usepackage{xcolor}
\usepackage{booktabs}
\usepackage{hyperref}
\usepackage{url}

% ---- listings configuration ----
\lstset{
  language=Python,
  frame=tb,
  numbers=left,
  numberstyle=\tiny\color{gray},
  backgroundcolor=\color{white},
  keywordstyle=\color{blue},
  stringstyle=\color{orange},
  commentstyle=\color{green},
  breaklines=true
}

% ---- minted style ----
\usemintedstyle{monokai}

\graphicspath{{./figures/}}

\title{Code Optimization Algorithms: A Comparative Study}

\author{Author Name}

\begin{document}

\maketitle

\begin{abstract}
We present a comparative study of optimization algorithms for code generation
and compilation. Our analysis covers gradient-based and gradient-free methods,
evaluating their effectiveness on various computational benchmarks.
\end{abstract}

\section{Introduction}

Algorithm efficiency is crucial in computational applications. This paper
compares state-of-the-art optimization algorithms for code generation tasks.

\section{Union-Find Algorithm}

The Union-Find data structure is fundamental in graph algorithms.

% ---- Algorithm environment ----
\begin{algorithm}[p]
  \caption{Union-Find with Path Compression}
  \begin{algorithmic}[1]
    \Function{MakeSet}{$x$}
      \State $\text{parent}[x] \gets x$
      \State $\text{rank}[x] \gets 0$
    \EndFunction
    \Statex
    \Function{Find}{$x$}
      \If{$\text{parent}[x] \neq x$}
        \State $\text{parent}[x] \gets \Call{Find}{\text{parent}[x]}$
      \EndIf
      \State \Return $\text{parent}[x]$
    \EndFunction
    \Statex
    \Function{Union}{$x, y$}
      \State $p \gets \Call{Find}{x}$
      \State $q \gets \Call{Find}{y}$
      \If{$p = q$}
        \State \Return
      \EndIf
      \If{$\text{rank}[p] < \text{rank}[q]$}
        \State $\text{parent}[p] \gets q$
      \ElsIf{$\text{rank}[p] > \text{rank}[q]$}
        \State $\text{parent}[q] \gets p$
      \Else
        \State $\text{parent}[q] \gets p$
        \State $\text{rank}[p] \gets \text{rank}[p] + 1$
      \EndIf
    \EndFunction
  \end{algorithmic}
\end{algorithm}

\section{Python Implementation}

The following Python code implements k-means clustering:

% ---- Python code block ----
\begin{lstlisting}[language=Python, caption={K-means clustering implementation}]
import numpy as np

def kmeans(X: np.ndarray, k: int, max_iter: int = 100) -> tuple:
    """
    K-means clustering algorithm.
    
    Args:
        X: Input data points
        k: Number of clusters
        max_iter: Maximum iterations
    
    Returns:
        Tuple of (centroids, labels)
    """
    centroids = X[:k]          # init: first k points
    labels = np.zeros(len(X), dtype=int)

    for _ in range(max_iter):
        # E-step: assign points to nearest centroid
        distances = np.linalg.norm(X[:, None] - centroids, axis=2)
        labels = np.argmin(distances, axis=1)

        # M-step: update centroids
        new_centroids = np.array([
            X[labels == i].mean(axis=0) if np.any(labels == i) else c
            for i, c in enumerate(centroids)
        ])

        if np.allclose(centroids, new_centroids):
            break
        centroids = new_centroids

    return centroids, labels
\end{lstlisting}

\section{Rust Implementation}

For performance-critical applications, we provide a Rust implementation:

% ---- Rust code block (minted) ----
\begin{minted}[caption={Graph adjacency list with Dijkstra}, frame=lines]{rust}
use std::collections::BinaryHeap;
use std::cmp::Reverse;

fn dijkstra(graph: &[Vec<(usize, i32)>], start: usize) -> Vec<i32> {
    let n = graph.len();
    let mut dist = vec![i32::MAX; n];
    dist[start] = 0;
    let mut pq = BinaryHeap::new();
    pq.push((Reverse(0), start));

    while let Some((Reverse(d), u)) = pq.pop() {
        if d > dist[u] { continue; }
        for &(v, w) in &graph[u] {
            if dist[u] + w < dist[v] {
                dist[v] = dist[u] + w;
                pq.push((Reverse(dist[v]), v));
            }
        }
    }
    dist
}
\end{minted}

\section{Performance Comparison}

\begin{table}[ht]
\centering
\caption{Execution time comparison (ms)}
\label{tab:perf}
\begin{tabular}{lccc}
  \toprule
  Algorithm & Python & Rust & Speedup \\
  \midrule
  K-means & 234.5 & 12.3 & 19.1x \\
  Dijkstra & 89.2 & 4.7 & 19.0x \\
  Sorting & 45.3 & 2.1 & 21.6x \\
  \bottomrule
\end{tabular}
\end{table}

\section{Conclusion}

Our comparative analysis reveals significant performance differences between
interpreted and compiled implementations.

\bibliographystyle{plain}
\bibliography{refs}

\end{document}
